\documentclass{article}
\usepackage[UTF8]{ctex}
\usepackage{amsmath,amssymb,amsfonts}
\usepackage{enumitem}
\usepackage{geometry}
\geometry{a4paper, margin=1in}
% === 错题本排版优化（不参与题目标识） ===
\setlength{\parindent}{0pt}
\setlength{\parskip}{0.35em plus 0.10em minus 0.08em}
\setlength{\emergencystretch}{3em}
\setlist{itemsep=0.25em, topsep=0.40em, parsep=0pt}
\allowdisplaybreaks[2]
% === 错题本排版优化结束 ===
\begin{document}
\section*{福州大学2025年数学分析真题}
\section*{117．福州大学2025年数学分析试卷}

\subsection{题目1}
1．（30 分）求极限．

\begin{itemize}
  \item[（1）] \(\lim _{x \rightarrow 0} \frac{\sqrt{1+x+x^{2}}+\sqrt{1+x-2 x^{2}}-2 \sqrt{1+x}}{e^{\cos x-1}-1}\) ．
  \item[（2）] \(\lim _{x \rightarrow \frac{\pi}{2}} \sum_{n=1}^{\infty} \frac{\sin (2 n+1) x}{n^{2}}\) ．
\end{itemize}

\subsection{题目2}
2．（20 分）叙述并证明数列的柯西收敛原理．

\subsection{题目3}
3．（20 分）设 \(f(x)\) 在 \([a,+\infty)\) 上连续，且 \(\lim _{x \rightarrow+\infty}[f(x)-2 \sqrt{x}]\) 存在，\(f(x)\) 在 \([a,+\infty)\) 上是否一致连续？如果是或否给出证明，不一定给出一正一反的例子。

\subsection{题目4}
4．（20 分）判别数项级数 \(\sum_{n=1}^{\infty} \frac{\sin n x}{\sqrt[4]{n}}\) 在 \((0, \pi)\) 上的敛散性，若收敛指出是绝对收敛还是条件收敛。

\subsection{题目5}
5．（20 分）证 \(\int_{-\infty}^{+\infty} e^{-\frac{x^{2}}{2}} \mathrm{~d} x=\sqrt{2 \pi}\) ，并用该结果计算 \(\int_{-\infty}^{+\infty} e^{-\frac{x^{2}}{2}+x} \mathrm{~d} x\) ．

\subsection{题目6}
6．（20 分）计算

\[
\iint_{\sum} \frac{2 x \mathrm{~d} y \mathrm{~d} z+(2+z)^{2} \mathrm{~d} x \mathrm{~d} y}{\sqrt{x^{2}+y^{2}+z^{2}}} .
\]

其中 \(\sum\) 是上半球面 \(z=\sqrt{4-x^{2}-y^{2}}\) ，方向取下侧．

\subsection{题目7}
7．（20 分）设 \(f(x)\) 在［0，1］上连续，定义 \(\alpha\)－阶积分算子 \((\alpha>0)\) ：

\[
I^{\alpha} f(x)=\frac{1}{\Gamma(\alpha)} \int_{0}^{x}(x-s)^{\alpha-1} f(s) \mathrm{d} s
\]

其中 \(\Gamma(\alpha)\) 为 Gamma 函数，利用积分换序原理证：当 \(\alpha, \beta>0\) 时，有

\[
I^{\alpha}\left(I^{\beta} f(x)\right)=I^{\alpha+\beta} f(x)=I^{\beta}\left(I^{\alpha} f(x)\right) .
\]
\end{document}
